\chapter{Knowledge and Skills Evaluation}
\label{chap:knowledge-skills-evaluation}

\noindent
This chapter evaluates the knowledge and skills developed throughout the project by comparing
them with the Program Learning Outcomes (PLOs) and the expected rubric levels established
by the curriculum. The evaluation is conducted from three perspectives: AI-based assessment,
student self-assessment, and advisor assessment. The purpose of this chapter is not only to
measure achievement, but also to provide reflective insight into how different evaluative
perspectives reveal strengths, gaps, and opportunities for further development.

\section{Comparison Between Program Targets and AI-Based Evaluation}
\label{sec:plo-ai-evaluation}

The first part focuses on comparing the program's expected learning outcomes with the
evaluation generated by AI. The AI should not only assign evaluation levels, but also provide
justifications explaining why each level was assigned.

\subsection{Radar Chart Comparing Program Targets and AI Evaluation}
\label{subsec:radar-ai}

This section includes a radar chart with two lines in the same figure: (1) the program's target
levels, and (2) the AI-based evaluation. The chart visually demonstrates the alignment or
discrepancy between the intended outcomes of the curriculum and the performance assessed
by AI.

\textit{[To be completed. Insert radar chart here.]}

% Example figure placeholder:
% \begin{figure}[H]
%     \centering
%     \includegraphics[width=0.75\textwidth]{figures/eval/radar_ai}
%     \caption{Radar chart comparing program targets and AI-based evaluation.}
%     \label{fig:radar-ai}
% \end{figure}

\subsection{Reflection on the AI Evaluation}
\label{subsec:reflection-ai}

This section discusses the student's opinions regarding the AI's assessment. The authors
reflect on whether the AI evaluation appears reasonable, insightful, too strict, too lenient, or
limited in any way, and provide supporting explanations.

\textit{[To be completed by the authors.]}

\section{Self-Evaluation by the Student}
\label{sec:self-evaluation}

The second part focuses on student self-evaluation. The student assesses their own knowledge
and skills and compares that self-assessment with the AI-generated evaluation. The discussion
explains why the two evaluations may differ.

\subsection{Radar Chart Comparing Program Targets, AI Evaluation, and Self-Evaluation}
\label{subsec:radar-self}

This section includes a radar chart with three lines in the same figure: (1) the program's target
levels, (2) the AI-based evaluation, and (3) the student's self-evaluation. This figure visually
illustrates similarities and differences among the three perspectives.

\textit{[To be completed. Insert radar chart here.]}

% Example figure placeholder:
% \begin{figure}[H]
%     \centering
%     \includegraphics[width=0.75\textwidth]{figures/eval/radar_self}
%     \caption{Radar chart comparing program targets, AI evaluation, and self-evaluation.}
%     \label{fig:radar-self}
% \end{figure}

\subsection{Reflection on Self-Evaluation Compared with AI Evaluation}
\label{subsec:reflection-self}

This section presents the student's interpretation of the self-evaluation in comparison with the
AI assessment. The student explains why they rated themselves at certain levels, why those
ratings may differ from the AI's conclusions, and what these differences reveal about
self-perception, evidence, and performance.

\textit{[To be completed by the authors.]}

\section{Advisor Evaluation}
\label{sec:advisor-evaluation}

The third part focuses on the evaluation conducted by the project advisor. The advisor assesses
the student's knowledge and skills and provides reasons for assigning each level.

\subsection{Radar Chart Comparing Program Targets, AI Evaluation, Self-Evaluation, and Advisor Evaluation}
\label{subsec:radar-advisor}

This section is completed only after the first two parts are finalized. The advisor reads the
previous sections and then conducts the third evaluation. A radar chart should be included with
four lines in the same figure: (1) the program's target levels, (2) the AI-based evaluation,
(3) the student's self-evaluation, and (4) the advisor's evaluation. This figure provides a
comprehensive comparison across all perspectives.

\textit{[To be completed by the advisor.]}

% Example figure placeholder:
% \begin{figure}[H]
%     \centering
%     \includegraphics[width=0.75\textwidth]{figures/eval/radar_advisor}
%     \caption{Radar chart comparing program targets, AI evaluation, self-evaluation, and advisor evaluation.}
%     \label{fig:radar-advisor}
% \end{figure}

\subsection{Advisor's Comments on the Evaluation}
\label{subsec:advisor-comments}

This section summarizes the advisor's opinions and explanations regarding the evaluation. The
advisor clarifies why each level was assigned and may comment on the student's strengths,
progress, limitations, and future development directions.

\textit{[To be completed by the advisor.]}

\subsection{Student Reflection on the Advisor's Evaluation}
\label{subsec:reflection-advisor}

This section presents the student's reflection on the advisor's assessment. The student may
discuss points of agreement or disagreement, lessons learned from the advisor's perspective,
and how the evaluation contributes to personal and professional development.

\textit{[To be completed by the authors.]}
